\section{已知内容防线：Hash Matching 的能力与边界}

当一个文件已经由可信机构验证并依法加入 reference set，最成熟的可扩展防线是 matching。它避免每次都让模型或人工重新理解同一内容，适合处理大规模重复传播；但它只能识别“与已知对象足够相同”的内容，不能发现所有新内容，也不能替代平台的合法性判断与报告流程。

\subsection{Cryptographic hash 与 perceptual hash}

\term{cryptographic hash}（密码学哈希）把字节序列映射为固定长度摘要，适合判断文件是否逐字节相同；轻微压缩或裁切就会改变摘要。\term{perceptual hash}（感知哈希）提取视觉或听觉结构，使经过常见变换的相似媒体仍可能落在邻近表示；它提高变体召回，却需要阈值、reference provenance 与误匹配控制。

\lecturefigure{03-hash-matching.png}{Hash matching 是对已知、经验证材料进行规模化阻断的最低防线。}{本地字幕 05:10--06:00；Thorn Safer Match 官方机制；概念重绘。}

\begin{knowledgebox}{读图：Match 不是“重新识别内容”}
平台先计算 incoming media 的 hash 或 fingerprint，再与受控 reference set 比对；命中后进入平台 policy、审核或报告路径。cryptographic match 通常具有更强的 byte identity 证据，perceptual match 则表达相似度。两者都需要记录 hash algorithm/version、reference source、threshold、match confidence 与后续人工动作，才能在系统升级后重放和审计。
\end{knowledgebox}

\begin{warningbox}{Hash database 不是普通 blocklist}
Reference hash 关联极高敏感度证据，必须限制导入来源、访问权限、导出能力和日志可见性。平台不能因为“只保存 hash”就忽视 membership inference、内部滥用、错误标注、跨租户泄露与删除义务。匹配服务应返回最小必要结果，而不是暴露 reference 内容或可枚举接口。
\end{warningbox}

\subsection{为什么已知匹配仍需人工与政策}

本节把 hash 信号接回响应链。相同文件出现在举报取证、新闻报道、研究、合法执法系统或用户上传环境中，其上下文和法定义务可能不同；不同司法辖区、服务类型和产品功能也有不同报告要求。因此工程团队应把“match”建模为高置信信号，而不是绕过所有 policy 的万能删除命令。

\begin{importantbox}{可审计匹配事件的最小字段}
一个 durable match event 至少包含：content identifier、hash algorithm/version、reference list/version、match distance、tenant/account、event time、policy version、review/action、reporting handoff 与 retry state。敏感字段应按最小权限分层存储，并为删除、legal hold 和 incident response 设计明确流程。
\end{importantbox}

\subsection{本章小结}

Hash matching 用已验证 reference set 换取速度与一致性，是 known-material defense 的基础。它的边界同样清楚：不能覆盖未知内容，不能推断行为人身份或意图，也不能免除上下文复核与法定程序。

